% Generated draft: requires mathematical and source review.
\documentclass[11pt]{article}
\usepackage[margin=25mm]{geometry}
\usepackage{fontspec}
\setmainfont{Latin Modern Roman}
\usepackage{amsmath,amssymb,mathtools,unicode-math}
\setmathfont{Latin Modern Math}
\usepackage{enumitem,needspace,xurl,hyperref}
\setlength{\parindent}{0pt}
\setlength{\parskip}{5pt}
\setlength{\emergencystretch}{3em}
\newcommand{\N}{\mathbb{N}}
\newcommand{\Z}{\mathbb{Z}}
\newcommand{\Q}{\mathbb{Q}}
\newcommand{\R}{\mathbb{R}}
\newcommand{\C}{\mathbb{C}}
\newcommand{\F}{\mathbb{F}}
\newcommand{\A}{\mathbb{A}}
\newcommand{\PP}{\mathbb{P}}
\newcommand{\E}{\mathbb{E}}
\newcommand{\eps}{\varepsilon}
\newcommand{\dd}{\,\mathrm{d}}
\begin{document}
\section*{UnsolvedMath ID: 30000645}
\textbf{Source problem code:} catalogue entry
\subsection{Definitions and mathematical statement}
\textbf{Top-Chern splitting question.} Let \(X=\operatorname{Spec}A\) be a smooth affine variety of dimension \(n\ge2\) over a field \(k\), and let \(P\) be a projective \(A\)-module of rank \(n\). If \(c_n(P)=0\) in \(CH_0(X)\), must there exist a projective module \(Q\) with\[P\cong A\oplus Q?\] The rank-one summand is free; the question does not require \(P\) or \(Q\) to be free. In the usual vector-bundle formulation, \(P\) is finitely generated of constant rank and corresponds to a rank-\(n\) bundle \(E\) on \(X\).\par For a smooth pure \(n\)-dimensional variety, its top Chern class in \(CH^n(X)\) is identified with a zero-cycle class in \(CH_0(X)\), the group of closed-point cycles modulo rational equivalence. The displayed condition concerns this ordinary Chow class, not vanishing in a refined Euler or oriented obstruction group. Splitting off \(A\) is equivalent to a unimodular element: there exist \(p_0\in P\) and \(\ell\in\operatorname{Hom}_A(P,A)\) with \(\ell(p_0)=1\). Geometrically this is the appropriate nowhere-vanishing algebraic section. Such a splitting forces the top Chern class to vanish; the source asks for the converse under its intended field and module hypotheses.\par The base field is unrestricted in the supplied wording. Algebraic closedness, characteristic, perfection, or infinitude are not stated and are not silently added. These are substantial scope choices in a splitting criterion, because the ordinary Chow class need not encode every field-dependent obstruction. Any comparison with a result requiring a special field or a refined obstruction theory must retain its hypotheses and cannot be treated as settling the literal all-field wording. The finite-generation, purity, and constant-rank conventions also require source confirmation if a broader module category is intended.\par The supplied triage checked 2026-08-21 labels the extracted question open and reports no located complete resolution, with the report as its sole reference. It gives no field-specific theorem or counterexample comparison. That dated status assertion is retained without independently certifying it for the literal unrestricted-field formulation.
\subsection{Short English statement}
Ask whether vanishing of the ordinary top Chern class forces a free rank-one direct summand in a rank-\(n\) projective module on a smooth affine \(n\)-fold. The field is unrestricted as written, and no algebraic-closedness or characteristic hypothesis is inserted. Refined obstruction vanishing would be a different premise.
\subsection{Sources}
\begin{itemize}\item Local UnsolvedMath registry, permanent ID 30000645, code catalogue entry; \emph{Splitting Projective Modules with Vanishing Top Chern Class}, Oberwolfach Reports Open Problems.\item \emph{Affine Algebraic Geometry (2007)}, Oberwolfach Reports, DOI 10.4171/owr/2007/01, pp. 5--82. \url{https://doi.org/10.4171/owr/2007/01}. Proposed year 2007; no proposer named.\item The same full report is the sole supplied literature source; no separate authoritative field-specific splitting theorem is supplied.\item Literature triage checked 2026-08-21, status \texttt{open}; record updated 2026-09-16. The \texttt{corrected\_verified} field records notation normalization and preservation of hypotheses. Verify dated claims independently.\end{itemize}
\end{document}
